% ====================================================================
%  Chapter 3 — Methods
% --------------------------------------------------------------------
%  Rewritten 2026-07-03 to describe the experiments actually run
%  (benchmark v0 + zero-shot baselines + training-free backend
%  adaptation + WavLM PEFT). Sources of record:
%    research_logs/2026-07-03-sl-benchmark-v0-design.md
%    research_logs/2026-07-03-zeroshot-baseline-table.md
%    research_logs/2026-07-03-p2-backend-adaptation.md
%    research_logs/2026-07-03-p3-peft-finetuning.md
%  The v1 (SLCeleb) protocol will be added once the corpus is staged.
% ====================================================================

\chapter{Methods}
\label{ch:methods}

\section{Corpus and benchmark construction}
\label{ssec:corpus}

We evaluate on a two-tier benchmark. \textbf{Benchmark v0}, used for
all experiments in this thesis, is built from two public read-speech
corpora: OpenSLR SLR52 \cite{kjartansson2018slr52} contributes 478
Sinhala speakers ($\approx$185k crowdsourced utterances, 16\,kHz FLAC,
minimum 10 utterances per speaker), and OpenSLR SLR65
\cite{he2020opensource} contributes 49 Tamil speakers (4{,}286 studio
utterances, 48\,kHz resampled at load). Speaker identities come from
the corpora's published metadata; gender labels exist only for SLR65.
A deterministic ingestion script (\texttt{tools/ingest\_openslr.py})
lays the corpus out as
\texttt{<lang>/<speaker>/<utterance>} symlinks and emits a manifest
and speaker-metadata table. \textbf{Benchmark v1}, pending data
staging, replaces the evaluation side with SLCeleb
\cite{slceleb2025}, our group's 280-speaker in-the-wild
Sinhala/Tamil celebrity corpus, which adds multi-genre audio,
multi-session speakers, and bilingual speakers for true
cross-lingual trials.

\section{Trial construction}
\label{ssec:trials}

Utterances are split 80/20 per speaker into train and test partitions
(seed 42, \texttt{tools/sl\_dataprep.py}). From the test partition we
sample, per language, 3{,}000 same-speaker target pairs and 9{,}000
same-language impostor pairs (1:3 ratio), with unordered-pair
deduplication and attempt-bounded sampling. Impostor pairs are
same-gender whenever both speakers' genders are known; this holds for
all Tamil trials but cannot be enforced for Sinhala (SLR52 has no
gender metadata). Cross-session targets are enforced when session
keys are derivable; the OpenSLR layouts carry no session metadata, so
v0 targets are not session-controlled. Consequently v0 absolute error
rates are optimistic and are valid for \emph{relative} comparison
between systems, not for comparison against in-the-wild benchmarks;
we state this caveat wherever v0 numbers are reported. The
cross-lingual trial list is empty at v0 (the OpenSLR corpora contain
no bilingual speakers) and becomes meaningful at v1.

\section{Metrics}
\label{ssec:metrics}

We report equal error rate computed as the mean of the false-positive
and false-negative rates at the crossing point (EER$_{\mathrm{avg}}$),
and minimum detection cost (minDCF) at two operating points,
$p_{\mathrm{target}} \in \{0.01, 0.05\}$ with
$C_{\mathrm{miss}}=C_{\mathrm{fa}}=1$, following NIST and VoxCeleb
practice. Calibration quality is measured with $C_{\mathrm{llr}}$ and
actual DCF after logistic calibration. Unlabeled trials are treated
as a hard error by the evaluation harness.

\section{Experiment suite A: zero-shot baselines}
\label{ssec:zeroshot}

Four released, VoxCeleb-trained checkpoints are evaluated with no
adaptation of any kind: ECAPA-TDNN \cite{desplanques2020ecapa} as
distributed by SpeechBrain \cite{ravanelli2021speechbrain}
($\approx$20M parameters), and ReDimNet-B1/B2/B6
\cite{yakovlev2024redimnet} (2.2M/4.7M/15M). Scoring is cosine
similarity on L2-normalised full-utterance embeddings, with no score
normalisation and no calibration; this table is the raw-transfer
reference that all adaptation levers are measured against.

\section{Experiment suite B: training-free backend adaptation}
\label{ssec:backend}

On top of cached suite-A embeddings we ablate three eval-time levers,
none of which updates the encoder:
(i)~\textbf{AS-Norm} \cite{matejka2017analysis}: adaptive symmetric
score normalisation with a 527-utterance in-language cohort (one
utterance per training-pool speaker), top-$K{=}300$;
(ii)~\textbf{PLDA} \cite{sizov2014unifying,kenny2010}: a
two-covariance model ($d{=}150$) trained on 10{,}480 utterances
($\leq$20 per speaker) from the training partition, preceded by
centering, LDA, and length normalisation;
(iii)~\textbf{language-aware logistic calibration}: an affine
score-to-LLR map fitted per language on half the trials and evaluated
on the other half, plus the cross-language transplant of each
calibrator to quantify the cross-lingual score shift
\cite{thienpondt2022calibration}. The grid crosses
$\{$cosine, PLDA$\}\times\{$raw, AS-Norm$\}$ for three encoders,
answering RQ4 directly. The PLDA fit necessarily sees the trial
speakers (utterance-disjoint, standard for domain-adaptation PLDA);
we flag this closed-set property explicitly.

\section{Experiment suite C: parameter-efficient fine-tuning}
\label{ssec:peft}

Suite C fine-tunes WavLM-base+ \cite{chen2022wavlm} (94.4M
parameters, English-pretrained) for Sinhala/Tamil speaker
verification. The 13 hidden states are combined by a learnable
softmax-weighted sum, pooled by attentive statistics pooling (128-d
bottleneck), and projected to a 192-d embedding trained with
AAM-Softmax \cite{deng2019arcface} (margin 0.2, scale 30) over the
527 training speakers. The CNN feature extractor stays frozen in all
arms. Four fine-tuning policies are compared:

\begin{description}
  \item[frozen] encoder frozen; only the pooling/embedding head and
  layer weights train ($\approx$0.5\% of parameters);
  \item[LoRA] rank-8 adapters ($\alpha{=}16$) \cite{hu2022lora} on the
  query/key/value/output projections of every attention block, plus
  the head (1.08M trainable, 1.1\%);
  \item[full] all encoder parameters unfrozen (encoder LR $10^{-5}$,
  head LR $10^{-3}$);
  \item[LLRD] full fine-tuning with layer-wise learning-rate decay
  ($\gamma{=}0.9$ per layer from the top).
\end{description}

Training uses the v0 train split capped at 60 utterances per speaker
(31{,}049 utterances), random 2.0\,s crops, batch 64, AdamW (weight
decay $2{\times}10^{-5}$), exponential LR decay ($\gamma{=}0.95$),
mixed precision, gradient clipping at 5.0, and 15 epochs. Model
selection uses per-epoch pooled EER on stratified 2{,}000-pair
subsamples; the selected checkpoint is finally evaluated on the full
12{,}000-pair lists, making suite-C numbers directly comparable to
suites A and B. Because the evaluation speakers are the training
speakers (utterance-disjoint split), suite C measures
\emph{representation adaptation} rather than open-set
generalisation; the v1 rerun provides the open-set condition.
Finally, the suite-B backend stack is re-applied to the fine-tuned
embeddings (suite B$\times$C composition) to test whether backend
adaptation retains value once the encoder is in-domain
\cite{wang2022plda}.

\section{Reproducibility}
\label{ssec:repro}

Trial lists, cohort lists, and PLDA training lists are reproducible
from seed 42. The LoRA and full arms of suite C are replicated with
three seeds (42/43/44) and reported as mean$\pm$std; the frozen and
LLRD arms use seed 42 only, as both sit far from the decision
boundary between policies. All runs, configurations, per-epoch
histories, and raw result JSONs are archived with the research logs
listed in the chapter header.
